\section{Discussion}

The results suggest a separation of roles. VQA confidence is the best available signal for ranking answer risk, while defect diagnosis is useful for explanation and recovery. This is plausible: a VQA model's confidence may already encode image quality, question difficulty, answer ambiguity, and model uncertainty in a single scalar. Image-quality defects are semantically meaningful, but they do not fully order examples by VQA correctness.

This negative result should not be hidden. In assistive settings, an honest reliability layer should distinguish between knowing that an answer is risky and knowing why the interaction failed. A high-confidence answer may be returned directly. A low-confidence answer may be refused. When refusing, the defect head can explain the likely visual cause and suggest a retake action. This design avoids overstating what defect labels can do for calibration while preserving their user-facing value.

The high false refilm rates indicate that recovery policies should be cost-sensitive. Asking a blind or low-vision user to retake an image has a real usability cost. Future systems should tune recovery thresholds to balance helpful retake guidance against unnecessary burden.
